\section{Introduction and quick start}\label{sec:introduction}

\CJKQC{} is a command-line program for turning folders of
quantum-chemistry output into a compact database, and then for inspecting,
filtering, and post-processing that database. It is a spiritual successor to JKQC from
the JKCS2.1 framework. The practical workflow is always the same:

\begin{enumerate}
  \item \textbf{Collect} program outputs into a database file, with one
        record per line.
  \item \textbf{Inspect and filter} the records: print columns, select
        compositions, apply energy or frequency cutoffs, remove duplicates,
        and sort.
  \item \textbf{Post-process} the survivors: derive thermochemistry,
        attach single-point corrections, compute formation energies, and
        print solvation distributions.
\end{enumerate}

Nothing derived is stored. A database holds only what the quantum-chemistry
program actually printed (energies, geometries, frequencies, and so on).
Zero-point energies, Gibbs free energies, corrections, formation energies,
populations, and distributions are recomputed on demand from those
stored values under the settings of the current run. This is why changing the
temperature, a frequency scale factor, or an energy unit never invalidates a
database.

\subsection{What can be read}

\CJKQC{} chooses how to read each file by looking at its \emph{content}, never
at its name or extension. The supported formats are:

\begin{center}
\begin{tabular}{@{}ll@{}}
\toprule
Program & Typical input \\
\midrule
Gaussian & \texttt{.log}, \texttt{.out} files with a Gaussian banner \\
ORCA     & \texttt{.out} files with ORCA's text banner \\
xtb      & log files with an \texttt{xtb version} banner \\
MRCC     & output files containing \texttt{MRCC program system} \\
plain xyz & an atom count, a comment line, and atom rows \\
\bottomrule
\end{tabular}
\end{center}

One file becomes exactly one record. For a file that chains several jobs,
the record describes the \emph{final} job (the last SCF energy, the final
geometry and frequencies, the final property block). For a multi-frame xyz
file, it describes the last complete frame. Failed and truncated jobs are
stored as records with status \fld{failed} or \fld{truncated}; they do not
abort a collection run. A file that matches no supported format, or an empty file,
\emph{does} abort a build run, because it is not a quantum-chemistry job at
all.

\subsection{Command anatomy}

Every invocation has the same shape:

\begin{verbatim}
CJKQC <inputs>... [filtering flags] [print flags]
      [settings flags] [-o out.jsonl]
\end{verbatim}

The positional \texttt{<inputs>} are either quantum-chemistry output files or
existing databases. \CJKQC{} classifies them by content:

\begin{itemize}
  \item \textbf{all inputs are databases} \(\rightarrow\) \textbf{merge mode}:
        the databases are combined, and records are copied exactly as stored;
  \item \textbf{all inputs are program outputs} \(\rightarrow\) \textbf{build mode}:
        each output is read into a record, compositions are resolved, and a
        new database is written;
  \item \textbf{a mixture of both kinds} is a usage error (exit code 2).
\end{itemize}

The file extension is irrelevant: a database named \texttt{weird.dat} still
merges if its first non-blank line begins with \texttt{\{}, and a Gaussian
log named \texttt{weird.xyz} is still read as Gaussian output.

\subsection{Three families of flags}

Flags fall into three groups, and the distinction matters for what a run is
allowed to do.

\begin{description}
  \item[Manipulation flags] change which records survive, or add corrected
        values to them:
        \flg{--extract}, \flg{--nextract}, the cutoff flags, \flg{--uniq},
        \flg{--reacted}, \flg{--reacted2}, \flg{--sort}, \flg{--reverse},
        \flg{--select}, \flg{--glob}, \flg{--shuffle}, \flg{--seed},
        \flg{--noname}, \flg{--add-correction}, and \flg{--extrapolate}.
  \item[Print flags] ask for output columns or sections: \flg{--filename},
        \flg{--el}, \flg{--g}, \flg{--pop}, \flg{--formation},
        \flg{--hydration}, \flg{--movie}, and the rest.
  \item[Settings flags] configure derivations or presentation but do not
        count as a manipulation or a print: \flg{--temp}, \flg{--fc},
        \flg{--fscale}, \flg{--makereal}, \flg{--ignoreimag}, \flg{--unit},
        \flg{--conc}, \flg{--cnt}, \flg{--bavg}, \flg{--rh}, and
        \flg{--psolvent}.
\end{description}

A run must have a \emph{visible outcome}. With neither manipulation nor print
flags, \flg{--out} is required and the run behaves exactly like a collection
run. A run with a manipulation flag but neither \flg{--out} nor a print flag
is a usage error, because the manipulated records would be stored nowhere and
nothing would be printed.

\begin{note}
A settings flag on its own is not a visible outcome. For example,
\texttt{CJKQC db.jsonl} \flg{--temp} \texttt{350} is a usage error; add a print flag or
\flg{--out}.
\end{note}

\subsection{Overwriting, summaries, and exit codes}

A file that already exists is never replaced silently. Immediately before the
old file is replaced, \CJKQC{} asks:

\begin{verbatim}
db.jsonl already exists; overwrite? [y/N]
\end{verbatim}

Only \texttt{y} or \texttt{yes} (any casing, surrounding whitespace ignored)
proceeds. When standard input is not an interactive terminal, the question is
not asked and the replacement is refused with a message telling you to use
\flg{--fout}. \flg{--fout} forces the replacement without asking.

The completion line printed by the program tells you what happened:

\begin{verbatim}
collected 120 records into db.jsonl (2 failed, 1 truncated)
merged 40 records into merged.jsonl
filtered 17 records into kept.jsonl
wrote 17 frames into movie.xyz
\end{verbatim}

Exit codes are:

\begin{center}
\begin{tabular}{@{}cl@{}}
\toprule
Code & Meaning \\
\midrule
0 & success, nothing skipped \\
1 & fatal error (nothing written) or a merge that skipped corrupt lines \\
2 & usage error (bad flag, missing \flg{--out}, mixed inputs) \\
\bottomrule
\end{tabular}
\end{center}

Failed and truncated records from a recognized format do not change the exit
code: they are data, not failures of the run.

\subsection{Quick start}

Collect every Gaussian and ORCA output in the current directory into a new
database, then print the names, electronic energies, and Gibbs free energies
of the structures within 10 kcal/mol of the lowest electronic energy of their
composition group:

\begin{verbatim}
# 1. Build a database (the shell expands *.out)
CJKQC *.out -o all.jsonl

# 2. Print a table
CJKQC all.jsonl --filename --el --g

# 3. Filter, print, and store the survivors
CJKQC all.jsonl --cutr el 10 --filename --el --g -o near-minima.jsonl
\end{verbatim}

The next chapters explain each step in detail, beginning with how data enters
the database and how to read the resulting database file.
